\chapter{Lessons Learned}

Booting the asynchronous kernel under QEMU exposed integration faults that unit
tests and static checks had missed. This appendix records the most instructive
ones.

\section{Bugs found on hardware}

\subsection{Panic handler recursed through heap-dependent logging}

\textbf{Symptom:} any fault during early boot silently triple-faulted instead of
printing a panic message. \textbf{Root cause:} the panic handler used the normal
logging path, which allocated from the kernel heap. A fault before the heap was
ready caused the panic handler to fault, re-entering itself, recursing until
stack overflow and triple fault. \textbf{Fix:} the panic handler now uses
\texttt{early\_logln!} (direct serial writes, no heap dependency). This unmasked
every subsequent bug.

\subsection{Blocked threads lost their execution context}

\textbf{Symptom:} a thread that blocked in \texttt{wait()} restarted from its
entry point when woken. \textbf{Root cause:} \texttt{block\_thread} cleared the
scheduler's \texttt{current\_handle} before \texttt{yield\_lp} saved the
thread's context, so the save was skipped. \textbf{Fix:} a self-blocking thread
stays \texttt{current\_handle} until the yield saves its context.

\subsection{RwLock held-across-call deadlocks}

\textbf{Symptom:} \texttt{sleep()} worked with logging but hung without it
(a heisenbug). \textbf{Root cause:} an \texttt{if let} binding extended an
RwLock read guard's lifetime across a call that tried to acquire the write lock.
\textbf{Fix:} bind the value out of the scrutinee first, so the temporary guard
drops before the re-locking call.

\subsection{Quantum timer grew without bound}

\textbf{Symptom:} the per-LP timer queue contained thousands of quantum events
after extended operation. \textbf{Root cause:} manual yields enqueued a new
quantum event each time. \textbf{Fix:} an \texttt{armed} flag keeps exactly one
quantum event in flight.

\subsection{A thread cannot free its own kernel stack}

\textbf{Symptom:} a returning kernel thread hung during teardown. \textbf{Root
cause:} \texttt{ThreadContext::Drop} deallocated the stack while the thread was
still executing on it. \textbf{Fix:} a deferred reaper moves dead threads to a
zombie list and drops them from the next thread's context after the switch.

\subsection{LA57 paging-mode mismatch}

\textbf{Symptom:} a \texttt{\#GP} with a non-canonical address on x86-64.
\textbf{Root cause:} the address width was derived from CPUID (57-bit capability)
instead of \texttt{CR4.LA57} (the active mode, 48-bit under Limine). \textbf{Fix:}
derive from the active paging mode.

\subsection{ASID-0 collision}

\textbf{Symptom:} the first user thread faulted at its entry point. \textbf{Root
cause:} the first user address space received ASID 0, which equals
\texttt{KERNEL\_ASID}, so the thread was created as a kernel thread (EL1),
where PXN forbade execution of user code. \textbf{Fix:} pre-populate the table
with the kernel AS at index 0.

\subsection{x86-64 trampoline halted on thread return}

\textbf{Symptom:} the async demo's worker completed but the coordinator never
resumed. \textbf{Root cause:} the x86-64 trampoline entered a \texttt{hlt} loop
when a thread returned, instead of calling \texttt{abort()}. \textbf{Fix:} make
the x86-64 trampoline call \texttt{abort} on return.

\subsection{Hypervisors that strip TTBR0 ASIDs break TLB isolation}

\textbf{Symptom:} under Apple's Hypervisor.framework, EL0 tests failed with
apparently unrelated data aborts, wrong syscall results, and deadline panics;
TCG remained green. \textbf{Root cause:} HVF did not preserve the ASID bits in
\texttt{TTBR0\_EL1}. Because context switches relied on those tags instead of a
full TLB flush, low user virtual addresses could resolve through another domain's
stale translation. \textbf{Fix:} \texttt{hvf\_compat} performs a whole-TLB flush
on each context switch and obtains syscall identity from trusted per-LP state.
\textbf{Lesson:} revalidate translation-tag assumptions on every hypervisor. See
\emph{AArch64 Port Status} for the detailed diagnosis.

\section{Engineering lessons}

\begin{enumerate}
    \item \textbf{Test transitions, not only interfaces.} The difficult failures
        occurred where block, yield, wake, context save, and reclamation met.
        Tests must cross those boundaries.

    \item \textbf{Keep failure paths independent.} Panic reporting must work
        before the heap and scheduler; reclamation must run from a live stack.

    \item \textbf{Validate platform assumptions at runtime.} CPU capability does
        not imply active paging mode, and a hypervisor may not preserve hardware
        translation tags.

    \item \textbf{Treat lock lifetime as part of control flow.} Bind values out of
        guards before calling code that may acquire the same lock.
\end{enumerate}

\endinput
